% Partie : sets-functions-relations
% Chapitre : size-of-sets
% Section : enumerations

\documentclass[../../../include/open-logic-section]{subfiles}

\begin{document}

\olfileid[fr]{sfr}{siz}{enm}

\olsection{Énumérations et ensembles \usetoken{p}{enumerable}}

\begin{editorial}
Cette section définit les énumérations d’ensembles comme des surjections
de $\PosInt$ sur ces ensembles. Elle procède progressivement, pour
les lecteurs ayant peu de connaissances mathématiques préalables.
La \olref[enm-alt]{sec} propose une présentation plus concise, fondée
sur une autre définition : une bijection avec $\Nat$ ou avec un segment
initial. Dans cette édition, \emph{au plus dénombrable} inclut les
ensembles finis et l’ensemble vide.\footnote{Convention éditoriale :
le mot anglais \emph{enumerable} est employé ici sans condition de
calculabilité. Une énumération n’est pas nécessairement donnée par
un algorithme ; il ne s’agit pas d’énumérabilité calculable.}
\end{editorial}

\begin{explain}
Nous avons déjà donné des exemples d’ensembles en dressant la liste
de leurs !!{element}s. Examinons plus généralement comment et dans
quels cas on peut dresser la liste des !!{element}s d’un ensemble,
même infini.
\end{explain}

\begin{defn}[Énumération : définition informelle]
Informellement, une \emph{énumération} d’un ensemble~$A$ est une liste,
éventuellement infinie, d’!!{element}s de~$A$, dans laquelle chaque
!!{element} de $A$ apparaît à une position finie. Si $A$ possède
une énumération, on dit qu’il est \emph{!!{enumerable}}.
\end{defn}

\begin{explain}
Précisons quelques points sur les énumérations :
\begin{enumerate}
\item Sauf pour la liste vide, admise ci-dessous, les listes considérées
  ont un premier !!{element}, et chaque !!{element} autre que le premier
  possède un unique !!{element} immédiatement avant lui. Nous exigeons
  en outre qu’il n’y ait qu’un nombre fini d’éléments entre le premier
  !!{element} et tout autre !!{element} de la liste. Ainsi, chaque
  !!{element} de l’énumération occupe une position finie : le premier
  est en position~$1$, le deuxième en position~$2$, etc.\footnote{Précision
  éditoriale : l’original présente ces deux conditions comme équivalentes.
  L’existence d’un prédécesseur immédiat ne suffit pas, à elle seule,
  à assurer que chaque terme occupe une position finie.}
\item Un même ensemble~$A$ peut avoir plusieurs énumérations, où les
  !!{element}s apparaissent dans des ordres différents : $4$, $1$,
  $25$, $16$,~$9$ énumère les cinq premiers carrés strictement positifs
  aussi bien que $1$, $4$, $9$, $16$,~$25$.
\item Une liste avec répétitions reste une énumération : $1$, $1$, $2$,
  $2$, $3$, $3$,~\dots{} énumère le même ensemble que $1$, $2$,
  $3$,~\dots
\item L’ordre et les répétitions \emph{comptent}, en revanche, lorsqu’on
  précise une énumération particulière. On peut énumérer les entiers
  strictement positifs en commençant par $1$, $2$, $3$, $1$, \dots{},
  mais la règle est plus facile à reconnaître dans l’énumération
  habituelle $1$, $2$, $3$, $4$,~\dots
\item Une énumération non vide doit avoir un début : \dots, $3$, $2$,
  $1$ n’est pas une énumération des entiers strictement positifs, car
  cette liste n’a pas de premier !!{element}. Pour comprendre le lien
  avec la définition informelle, demandez-vous à quelle position
  de cette liste se trouverait le nombre 76.
\item La liste suivante n’énumère pas non plus les entiers strictement
  positifs : $1$, $3$, $5$, \dots, $2$, $4$, $6$, \dots\@ Les nombres
  pairs y seraient placés aux positions $\infty + 1$, $\infty + 2$,
  $\infty + 3$, et non à des positions finies.
\item L’ensemble vide est au plus dénombrable : la liste vide l’énumère !
\end{enumerate}
\end{explain}

\begin{prop}
Si $A$ possède une énumération, il en possède une sans répétitions.
\end{prop}

\begin{proof}
Si $A$ est vide, sa liste vide ne comporte aucune répétition.
Sinon, supposons que $A$ possède une énumération $x_1$, $x_2$, \dots{}, dont
chaque terme $x_i$ est un !!{element} de~$A$. On peut supprimer les
répétitions en retirant les termes déjà rencontrés : on conserve $x_i$
s’il ne figure pas parmi $x_1$, \dots, $x_{i-1}$, et on le retire
s’il apparaît déjà parmi $x_1$, \dots,~$x_{i-1}$.
\end{proof}

Ce dernier raisonnement montre que, pour mieux comprendre les
énumérations et les ensembles !!{enumerable}s, puis démontrer des
résultats à leur sujet, il nous faut une définition plus précise.
La voici.

\begin{defn}[Énumération : définition formelle]
Une \emph{énumération} d’un ensemble $A \neq \emptyset$ est une
fonction !!{surjective} $f \colon \PosInt \to A$.
\end{defn}

\begin{explain}
Vérifions que cette définition formelle correspond à la définition
informelle par une liste éventuellement infinie. Toute fonction
!!{surjective} de $\PosInt$ dans un ensemble~$A$ énumère~$A$ : elle
détermine la liste $f(1)$, $f(2)$, $f(3)$, \dots. La surjectivité de $f$
assure que chaque !!{element} de~$A$ est la valeur de~$f(n)$ pour
un certain~$n \in \PosInt$. Chaque !!{element} apparaît donc à une
position finie. La fonction n’étant pas nécessairement !!{injective},
la liste peut comporter des répétitions ; nous avons vu que cela
est permis.

Réciproquement, à partir d’une liste qui énumère tous les !!{element}s
de~$A$, on peut définir une fonction !!{surjective} $f\colon \PosInt \to A$
en prenant pour $f(n)$ le $n$-ième !!{element} de la liste, ou son dernier
!!{element} si elle n’a pas de $n$-ième !!{element}. Cette construction
échoue seulement lorsque $A$ est vide, et que la liste est donc vide.
Ainsi, toute liste non vide qui énumère~$A$ détermine une fonction
!!{surjective} $f\colon \PosInt \to A$.
\end{explain}

\begin{defn}
  \ollabel{defn:enumerable}
Un ensemble~$A$ est !!{enumerable} si et seulement s’il est vide
ou possède une énumération.
\end{defn}

\begin{ex}
Pour énumérer les entiers strictement positifs ($\PosInt$), il suffit
de la fonction identité $f(n) = n$. La fonction $g(n) = n - 1$
énumère les entiers naturels $\Nat$.
\end{ex}

\begin{ex}
Les fonctions $f\colon \PosInt \to \PosInt$ et
$g \colon \PosInt \to \PosInt$ définies par
\begin{align*}
f(n) & = 2n \text{ et}\\
g(n) & = 2n - 1
\end{align*}
énumèrent respectivement les entiers pairs strictement positifs et
les entiers impairs strictement positifs. Cependant, aucune des deux
n’énumère $\PosInt$, car aucune n’est !!{surjective} sur cet ensemble.
\end{ex}

\begin{prob}
Définissez une énumération des carrés strictement positifs
$1$, $4$, $9$, $16$, \dots
\end{prob}

\begin{ex}
La fonction $f(n) = (-1)^{n} \lceil \frac{(n-1)}{2}\rceil$ énumère
l’ensemble des entiers relatifs~$\Int$. Ici, $\lceil x \rceil$ désigne
le \emph{plafond} de~$x$, c’est-à-dire le plus petit entier supérieur
ou égal à~$x$. Les valeurs de $f$ passent alternativement des entiers
positifs aux entiers négatifs :
\[
\begin{array}{c c c c c c c c}
f(1) & f(2) & f(3) & f(4) & f(5) & f(6) & f(7) & \dots \\ \\
- \lceil \tfrac{0}{2} \rceil & \lceil \tfrac{1}{2}\rceil & - \lceil \tfrac{2}{2} \rceil & \lceil \tfrac{3}{2} \rceil & - \lceil \tfrac{4}{2} \rceil  & \lceil \tfrac{5}{2}
\rceil & - \lceil \tfrac{6}{2} \rceil & \dots \\ \\
0 & 1 & -1 & 2 & -2 & 3 & -3 & \dots
\end{array}
\]
\footnote{Correction éditoriale : on rétablit la valeur $-3$ sous
$f(7)$ dans la dernière ligne du tableau ; elle manquait dans l’original.}
On peut également définir $f$ par cas :
\[
f(n) = \begin{cases}
  0 & \text{si $n = 1$}\\
  n/2 & \text{si $n$ est pair}\\
  -(n-1)/2 & \text{si $n$ est impair et $>1$}
  \end{cases}
\]
\end{ex}

\begin{prob}
Montrez que, si $A$ et $B$ sont !!{enumerable}s, $A \cup B$ l’est aussi.
Pour cela, supposez données des fonctions !!{surjective}s
$f\colon \PosInt \to A$ et $g\colon \PosInt \to B$. Définissez une
fonction~$h\colon \PosInt \to A \cup B$ et montrez qu’elle est
!!{surjective}. Traitez aussi les cas où $A$ ou~$B = \emptyset$.
\end{prob}

\begin{prob}
Montrez que, si $B \subseteq A$ et $A$ est !!{enumerable}, $B$ l’est
aussi. Supposez pour cela donnée une fonction !!{surjective}
$f\colon \PosInt \to A$. Définissez une fonction~$g\colon \PosInt \to B$
et montrez qu’elle est !!{surjective}. Que se passe-t-il si
$B = \emptyset$ ?
\end{prob}

\begin{prob}
Montrez par récurrence sur $n$ que, si $A_1$, $A_2$, \dots, $A_n$ sont
tous !!{enumerable}s, $A_1 \cup \dots \cup A_n$ l’est aussi.
Vous pouvez admettre que la réunion de deux ensembles $A$ et~$B$
!!{enumerable}s, $A \cup B$, est !!{enumerable}.
\end{prob}

Lorsqu’on dresse la liste des !!{element}s d’un ensemble, il paraît
peut-être plus naturel de commencer par le premier. Les mathématiciens
utilisent pourtant volontiers les entiers naturels~$\Nat$ pour
numéroter les objets : ils parlent du terme d’indice~$0$, puis des
termes d’indices~$1$, $2$, etc. On peut donc aussi définir une énumération
comme une fonction !!{surjective} de $\Nat$ dans~$A$.
Les deux définitions sont équivalentes.

\begin{prop}\ollabel{prop:enum-shift}
Il existe !!a{surjection} $f\colon \PosInt \to A$ si et seulement
s’il existe !!a{surjection} $g\colon \Nat \to A$.
\end{prop}

\begin{proof}
Étant donnée !!a{surjection} $f\colon \PosInt \to A$, posons
$g(n) = f(n+1)$ pour tout $n \in \Nat$. On vérifie aisément que
$g\colon \Nat \to A$ est !!{surjective}. Réciproquement, étant donnée
!!a{surjection} $g\colon \Nat \to A$, posons $f(n) = g(n-1)$.
\end{proof}

On obtient ainsi le résultat suivant :

\begin{cor}\ollabel{cor:enum-nat}
Un ensemble $A$ est !!{enumerable} si et seulement s’il est vide
ou s’il existe une fonction !!{surjective} $f\colon \Nat \to A$.
\end{cor}

Nous avons vu qu’on peut transformer une liste d’!!{element}s
d’un ensemble~$A$ en une liste sans répétitions. Cela vaut aussi pour
les énumérations au sens formel, mais l’énoncé et la démonstration
demandent davantage de soin. Toute fonction $f\colon \PosInt \to A$
doit être définie pour chaque $n \in \PosInt$. Si $A$ ne possède
qu’un nombre fini d’!!{element}s, une fonction définie sur tous les
!!{element}s de $\PosInt$ ne peut prendre pour valeurs tous les
!!{element}s de~$A$ sans jamais se répéter. Lorsque la liste sans
répétitions est finie, il faut donc choisir pour~$f$ un autre domaine,
qui ne possède qu’un nombre fini d’!!{element}s. L’absence de répétitions
signifie que $f$ est !!{injective}. Puisqu’elle est aussi !!{surjective},
nous cherchons !!a{bijection} entre $\{1, \dots, n\}$ ou $\PosInt$ et~$A$.

\begin{prop}\ollabel{prop:enum-bij}
Si $f\colon \PosInt \to A$ est !!{surjective}, c’est-à-dire si elle
énumère~$A$, il existe !!a{bijection} $g\colon Z \to A$, où $Z$ est
soit~$\PosInt$, soit $\{1, \dots, n\}$ pour un certain~$n \in \PosInt$.
\end{prop}

\begin{proof}
Définissons $g$ par récurrence. Posons $g(1) = f(1)$. Une fois la valeur
$g(i)$ définie, prenons pour $g(i+1)$ la première valeur de la liste $f(1)$,
$f(2)$, \dots{} qui ne figure pas déjà parmi $g(1)$, \dots, $g(i)$,
s’il en existe une. Si $A$ possède exactement $n$ !!{element}s,
toutes les valeurs $g(1)$, \dots, $g(n)$ sont définies, ce qui donne
une fonction $g\colon \{1, \dots, n\} \to A$. Si $A$ possède une
infinité d’!!{element}s, alors, pour tout $i$, l’énumération $f(1)$,
$f(2)$, \dots{} contient un !!{element} de~$A$ qui ne figure pas encore
parmi $g(1)$, \dots, $g(i)$. On obtient alors une fonction
$g\colon \PosInt \to A$.

La fonction $g$ est !!{surjective} : tout élément de~$A$ apparaît
parmi $f(1)$, $f(2)$, \dots{}, puisque $f$ est !!{surjective}, et finit
donc par être une valeur $g(i)$. Elle est aussi !!{injective} : s’il
existait $j < i$ tel que $g(j) = g(i)$, la valeur $g(i)$ figurerait déjà
parmi $g(1)$, \dots, $g(i-1)$, contrairement à la définition de~$g$.
\end{proof}

\begin{cor}\ollabel{cor:enum-nat-bij}
Un ensemble $A$ est !!{enumerable} si et seulement s’il est vide
ou s’il existe !!a{bijection} $f\colon N \to A$, où $N = \Nat$
ou $N = \{0, \dots, n\}$ pour un certain $n \in \Nat$.
\end{cor}

\begin{proof}
L’ensemble $A$ est !!{enumerable} si et seulement s’il est vide
ou s’il existe une fonction !!{surjective} $f\colon \PosInt \to A$.
D’après la \olref{prop:enum-bij}, cette dernière condition équivaut
à l’existence d’une fonction !!{bijective}~$f\colon Z \to A$, où
$Z = \PosInt$ ou $Z = \{1, \dots, n\}$ pour un certain $n \in \PosInt$.
Le même décalage d’indice que dans la preuve de la \olref{prop:enum-shift}
montre que cette condition équivaut à l’existence d’!!a{bijection}
$g\colon N \to A$, où $N = \Nat$ ou $N = \{0, \dots, n-1\}$.
\end{proof}

\begin{prob}
D’après la \olref[sfr][siz][enm]{defn:enumerable}, un ensemble $A$
est au plus dénombrable si et seulement si $A = \emptyset$ ou s’il
existe une fonction !!{surjective} $f\colon \PosInt \to A$.
On peut aussi définir les ensembles !!{enumerable}s par l’existence
d’une fonction !!{injective} $g\colon A \to \PosInt$. Montrez que
les deux définitions sont équivalentes : il existe une fonction
!!{injective} $g\colon A \to \PosInt$ si et seulement si $A = \emptyset$
ou s’il existe une fonction !!{surjective} $f\colon \PosInt \to A$.
\end{prob}

\end{document}
